\section{Conclusion}

Under a fixed test-time budget, theorem proving is not only a generation
problem but a decision about which computation to fund next. We formulated
this decision as multi-resource constrained metareasoning and presented
\system, a controller over generation, repair, retrieval, capability tests,
decomposition, representation change, and model escalation. A replayable
workspace exposes the dependencies, alternatives, diagnostics, and verified
artifacts needed to make those allocation decisions auditable.

The theory intentionally stops short of a global optimality claim. Hard-budget
selection is already a knapsack problem in the static independent case;
\(p_i/c_i\) optimally orders only a fixed feasible set, and cheap-first
escalation requires a feasible fallback reserve. Residual verifier progress is
an empirical approximation to the continuation value that becomes necessary
when proof actions change the state.

The evaluation is designed around the corresponding scientific claim. Strong
minimal refinement tests whether allocation gains cover structured overhead,
while matched structured contrasts separately test action-space expansion,
cost-aware selection, remaining-budget admission, and adaptive routing.
\todo{Close with the absolute accepted-rate difference, paired confidence
interval, unconditional cost difference, exact budget vector, and the task
regime in which the result holds.}
